% !PRINTED 310
\weberpage{310}

We wish first to apply the transformation equation \S\,72, (8)',
in which now, in agreement with the preceding section,
$v$ is put for $x$:
\weq{1}{v^6 + 10\,v^3 - \gamma_2 v + 5 = 0}
Here, if
\[
c \equiv 0 \ (\mathrm{mod}\ 12), \quad z \equiv c \ (\mathrm{mod}\ 5)
\]
then
\weq{2}{v_\infty = 5\left(\frac{\eta(5\omega)}{\eta\,\omega}\right)^2, \quad v_z = \left(\frac{\eta\left(\dfrac{\omega+c}{5}\right)}{\eta(\omega)}\right)^2.} % [sic] The source prints \eta\omega without function parentheses in the first denominator; the parallel second denominator prints \eta(\omega).
If we then take
\weq{3}{\begin{aligned}
\sqrt{5}\,w_z &= (v_\infty - v_z)(v_{z+1} - v_{z-1})(v_{z-2} - v_{z+2}) \\
&= (v_\infty - v_c)(v_{c+12} - v_{c-12})(v_{c+24} - v_{c-24}), \\
z &= 0, 1, 2, 3, 4, \quad c = 0, \pm 12, \pm 24,
\end{aligned}}
then, by the concluding remark of the preceding section, the
$w_0$, $w_1$, $w_2$, $w_3$, $w_4$ are the roots of an equation of degree 5,
whose coefficients are formed rationally from $\gamma_2$ and rational
numbers.

If, however, we observe the relations [\S\,54, (14)]:
\[
\gamma_2(\omega+1) = e^{-\frac{2\pi i}{3}} \gamma_2(\omega), \quad \gamma_2\left(-\frac{1}{\omega}\right) = \gamma_2(\omega),
\]
it follows from the form of (1) that under the interchange
$(\omega, \omega+1)$ the roots $v$, apart from a rearrangement, take on the
factor $e^{\frac{2\pi i}{3}}$, while under $\left(\omega, -\dfrac{1}{\omega}\right)$ they remain
unchanged.

The $w_z$ are therefore interchanged only among themselves, and
their symmetric functions remain \gesperrt{unchanged} and depend
therefore rationally only on the invariant $j(\omega)$. The coefficients
in the equation for $w$ can moreover become infinite for no finite $j(\omega)$,
and they are accordingly \gesperrt{polynomials}
in $j(\omega)$.

Further information about these coefficients results
from the expansion in ascending powers of $q$. By (2) the
expansions of the $v$ have the beginnings
\[
v_\infty = 5\,q^{\frac{2}{3}} + \cdots, \quad v_z = e^{\frac{c}{12}\frac{2\pi i}{5}} q^{-\frac{2}{15}} + \cdots,
\]
and from this it follows, with the use of the well-known equation
\[
4\sin\frac{2\pi}{5}\cdot\sin\frac{4\pi}{5} = \sqrt{5}
\]
%CONT>>
